tr\chapter{Problema de Investigación}
Se espera que la descripción del problema sea amplia y detallada, se sugiere como mínimo tres páginas - según guía.

La problemática inicia cuando el sujeto detecta una necesidad concreta, falta de conocimiento o una contradicción entre los enfoques disponibles.

El problema de investigación se refiere a al descripción amplia y detallada de una realidad en la cual diversas variables o factores se relacionan, hechos-causas y consecuencias, y se concluye con una pregunta de investigación, que en esencia es la síntesis de la situación descrita. Es la conclusión a la que se llega después de haber descrito la situación problemática - Más información ver guía.

\section{Problema General}
Se sugiere que sea el titulo de la tesis en forma de pregunta. Se hace una pregunta la cual es respondida por la tesis. ¿Qué es lo difícil?


\section{Problemas Específicos}
(Se sugiere que sean pocos pero sustanciales)
El estudio toma en cuenta los siguientes problemas específicos:
\begin{itemize}
	\item Primer problema especifico.
	\item Segundo problema especifico.
	\item Tercer problema especifico.	
\end{itemize}
